\begin{ExerciseBlock}
\ExerciseHeading{}{Verständnisfragen}{linear_kinematics_checks}
Entscheiden Sie, welche Aussagen wahr oder falsch sind, und begründen Sie Ihre Entscheidung.
\begin{enumerate}
\item Die SI-Einheit einer Distanz ist Kilometer (km).\TrueFalseAnswer{falsch}
\item Nur Elementarteilchen (Elektronen, ...) lassen sich näherungsweise als punktförmige Objekte beschreiben.\TrueFalseAnswer{falsch}
\item Die SI-Einheit der Zeit ist die Sekunde (s).\TrueFalseAnswer{wahr}
\item Um eine Geschwindigkeit von m/s in km/h umzurechnen, muss man mit dem Faktor $3,6$ multiplizieren.\TrueFalseAnswer{wahr}
\item Die Momentangeschwindigkeit $v(t_0)$ ist der Grenzwert der mittleren Geschwindigkeit für ein gegen null gehendes Zeitintervall um $t_0$.\TrueFalseAnswer{wahr}
\item Das bestimmte Integral ergibt den vorzeichenbehafteten Flächeninhalt zwischen Graph und Achse.\TrueFalseAnswer{wahr}
\item Aus der Momentangeschwindigkeit $v(t)$ erhält man die Ortsänderung durch zeitliche Integration.\TrueFalseAnswer{wahr}
\item Die Geschwindigkeitsfunktion $v(t)$ einer gleichförmig beschleunigten Bewegung ist eine Gerade.\TrueFalseAnswer{wahr}
\item Bei einer gleichförmigen Bewegung hat der Graph der Beschleunigung $a(t)$ eine von null verschiedene Steigung.\TrueFalseAnswer{falsch}
\item Wenn der freie Fall aus einer Höhe $s$ die Fallzeit $T$ andauert, so dauert der freie Fall aus der Höhe $2s$ auch $2T$ an.\TrueFalseAnswer{falsch}
\item Im Modell des freien Falls nahe der Erdoberfläche ist $g\approx\qty{9.81}{\metre\per\second\squared}$ konstant und zum Erdmittelpunkt gerichtet.\TrueFalseAnswer{wahr}
\item Ein Einheitencheck allein genügt als Nachweis dafür, dass eine Formel korrekt ist.\TrueFalseAnswer{falsch}
\item Für ein symmetrisches Intervall $[t_0 - \Delta t, t_0 + \Delta t]$ um $t_0$ entspricht die mittlere Geschwindigkeit $\bar{v}$ immer der Momentangeschwindigkeit $v(t_0)$.\TrueFalseAnswer{falsch}
\end{enumerate}
\end{ExerciseBlock}
